\rtmtight
\begin{tblr}{width=\textwidth,colspec={|Q[c,m,2.1cm]|X[l,m]|},hlines,vlines,hspan=minimal,rowsep=1.5pt,colsep=3pt,row{1}={bg=headgray,font=\bfseries}}
\SetCell{c,m}\hfil\logo\hfil & \textbf{POLITEKNIK NEGERI MALANG}\par\textbf{JURUSAN TEKNOLOGI INFORMASI}\par\textbf{PROGRAM STUDI: D4 TEKNIK INFORMATIKA} \\
\SetCell[c=2]{l} \textbf{RENCANA TUGAS MAHASISWA (RTM) DAN RUBRIK PENILAIAN} & \\
Nama Mata Kuliah & Manajemen Proyek \\
Kode Mata Kuliah & RTI253001 \quad \textbf{sks / Jam perminggu:} 2 SKS/3 jam \quad \textbf{Semester:} 3 \\
Dosen Pengampu & \begin{enumerate}\item Pramana Yoga Saputra, S.Kom., M.MT.\item Ahmadi Yuli Ananta, S.T., M.M.\item Erfan Rohadi, S.T., M.Eng., Ph.D.\end{enumerate} \\
\end{tblr}
\assessmentblock{Tugas Proyek Kelompok Manajemen Proyek}{Project Based Learning dan presentasi tugas}{SCPMK0801-01901--SCPMK0402-01904}{Mahasiswa dalam kelompok menerapkan knowledge area manajemen proyek (Tugas 1--11) pada satu proyek perangkat lunak sepanjang semester: pemilihan topik, integrasi, pemangku kepentingan, sumber daya, WBS, jadwal, biaya, kualitas, komunikasi, risiko, Agile, hingga dokumentasi penutupan.}{Membentuk kelompok, menentukan topik proyek, menerapkan setiap knowledge area pada proyek kelompok sesuai penugasan mingguan, mendokumentasikan hasil, dan mempresentasikan progres di kelas.}{Dokumen rencana proyek (proposal, WBS, timeline, anggaran), artefak penerapan knowledge area, dokumentasi penutupan proyek, dan bahan presentasi.}{Kualitas perencanaan proyek (integrasi, scope/WBS, jadwal, biaya) & 10 & Rencana proyek lengkap: topik jelas, WBS terstruktur, timeline dan anggaran realistis serta konsisten antarartefak. & Rencana proyek tersedia, tetapi sebagian artefak (WBS, timeline, atau anggaran) kurang rinci atau kurang konsisten. & Rencana proyek tidak lengkap atau tidak dapat dijadikan dasar eksekusi proyek. \\ \\
Penerapan sumber daya, kualitas, risiko, dan Agile pada eksekusi proyek & 8 & Knowledge area eksekusi diterapkan pada proyek kelompok dengan bukti artefak dan tindak lanjut yang jelas. & Sebagian knowledge area diterapkan, tetapi bukti penerapan atau tindak lanjutnya terbatas. & Penerapan knowledge area tidak terlihat pada proyek kelompok. \\ \\
Pengelolaan pemangku kepentingan dan kerja sama tim & 2 & Identifikasi dan pengelolaan pemangku kepentingan terdokumentasi; peran anggota tim jelas dan komunikasi efektif. & Pemangku kepentingan teridentifikasi, tetapi pengelolaan hubungan atau pembagian peran tim kurang jelas. & Tidak ada pengelolaan pemangku kepentingan dan kerja tim tidak terkoordinasi. \\ \\
Dokumentasi, evaluasi, dan simulasi penutupan proyek & 6 & Dokumentasi proyek sistematis, evaluasi dan lessons learned jelas, simulasi proyek berjalan sesuai rencana. & Dokumentasi tersedia, tetapi evaluasi atau lessons learned belum lengkap. & Dokumentasi proyek tidak sistematis atau tidak menggambarkan pelaksanaan proyek. \\}

\assessmentblock{Kuis 1 Konsep dan Pemangku Kepentingan Proyek}{Kuis (tes tulis/online)}{SCPMK0801-01901, SCPMK0302-01903}{Mahasiswa menjawab soal kuis materi pertemuan 1--3: konsep proyek dan manajemen proyek, manajemen integrasi, serta manajemen pemangku kepentingan.}{Mengerjakan soal tes tulis/online secara individu dalam waktu 1x50 menit sesuai materi pertemuan 1--3.}{Lembar jawaban tes tulis atau jawaban pada platform ujian online.}{Penguasaan konsep proyek, fase, knowledge area, dan integrasi proyek & 7 & Jawaban tepat menjelaskan konsep proyek, fase, 10 knowledge area, dan proses manajemen integrasi. & Sebagian jawaban tepat, tetapi ada konsep atau proses yang keliru atau tidak lengkap. & Jawaban banyak yang keliru atau tidak menunjukkan pemahaman konsep. \\ \\
Ketepatan penjelasan manajemen pemangku kepentingan & 3 & Identifikasi, perencanaan, dan pengelolaan pemangku kepentingan dijelaskan tepat pada kasus soal. & Penjelasan pemangku kepentingan cukup tepat, tetapi belum lengkap pada kasus soal. & Penjelasan pemangku kepentingan keliru atau tidak relevan dengan soal. \\}

\assessmentblock{UTS Perencanaan Proyek}{UTS (tes tulis/online)}{SCPMK0801-01901, SCPMK0803-01902, SCPMK0302-01903}{Mahasiswa menjawab soal UTS materi pertemuan 1--7: konsep dan integrasi proyek, pemangku kepentingan, sumber daya dan pengadaan, ruang lingkup (WBS), serta manajemen waktu.}{Mengerjakan soal tes tulis/online secara individu dalam waktu 1x50 menit sesuai materi pertemuan 1--7.}{Lembar jawaban tes tulis atau jawaban pada platform ujian online.}{Ketepatan perencanaan proyek: integrasi, scope/WBS, jadwal & 15 & Jawaban menunjukkan penguasaan perencanaan proyek: proses integrasi, penyusunan WBS, dan penjadwalan diterapkan tepat pada kasus. & Sebagian besar jawaban tepat, tetapi ada langkah perencanaan yang keliru atau tidak lengkap. & Jawaban tidak menunjukkan penguasaan proses perencanaan proyek. \\ \\
Ketepatan penjelasan sumber daya dan pengadaan proyek & 5 & Proses perencanaan, perolehan, dan pengendalian sumber daya serta pengadaan dijelaskan tepat. & Penjelasan sumber daya/pengadaan cukup tepat, tetapi belum lengkap. & Penjelasan sumber daya/pengadaan keliru atau tidak relevan. \\ \\
Ketepatan penjelasan pemangku kepentingan dan organisasi proyek & 5 & Struktur organisasi dan pengelolaan pemangku kepentingan dijelaskan tepat pada kasus soal. & Penjelasan cukup tepat, tetapi kurang lengkap pada kasus soal. & Penjelasan keliru atau tidak menjawab soal. \\}

\assessmentblock{Kuis 2 Biaya, Kualitas, Risiko, dan Komunikasi}{Kuis (tes tulis/online)}{SCPMK0803-01902, SCPMK0302-01903}{Mahasiswa menjawab soal kuis materi pertemuan 9--11: manajemen biaya, kualitas, komunikasi, dan risiko proyek.}{Mengerjakan soal tes tulis/online secara individu dalam waktu 1x50 menit sesuai materi pertemuan 9--11.}{Lembar jawaban tes tulis atau jawaban pada platform ujian online.}{Ketepatan penjelasan manajemen biaya, kualitas, dan risiko & 7 & Perkiraan biaya, penjaminan/pengendalian kualitas, serta identifikasi dan strategi risiko dijelaskan tepat pada kasus. & Sebagian jawaban tepat, tetapi ada proses yang keliru atau tidak lengkap. & Jawaban banyak yang keliru atau tidak menunjukkan pemahaman materi. \\ \\
Ketepatan penjelasan manajemen komunikasi proyek & 3 & Tujuan dan proses manajemen komunikasi proyek dijelaskan tepat pada kasus soal. & Penjelasan komunikasi proyek cukup tepat, tetapi belum lengkap. & Penjelasan komunikasi proyek keliru atau tidak relevan. \\}

\assessmentblock{UAS Manajemen Proyek}{UAS (tes tulis/online)}{SCPMK0801-01901--SCPMK0402-01904}{Mahasiswa menjawab soal UAS materi pertemuan 1--16 yang mencakup seluruh knowledge area, Agile, penutupan proyek, dan dokumentasi proyek.}{Mengerjakan soal tes tulis/online secara individu dalam waktu 1x50 menit sesuai materi pertemuan 1--16.}{Lembar jawaban tes tulis atau jawaban pada platform ujian online.}{Penguasaan komprehensif konsep dan perencanaan proyek & 8 & Jawaban mengaitkan konsep, fase, knowledge area, dan perencanaan proyek secara tepat pada kasus. & Sebagian besar jawaban tepat, tetapi keterkaitan antarkonsep belum lengkap. & Jawaban tidak menunjukkan penguasaan konsep dan perencanaan proyek. \\ \\
Ketepatan eksekusi, monitoring, Agile, dan penutupan proyek & 8 & Proses eksekusi-monitoring, prinsip Agile, dan penutupan proyek dijelaskan dan diterapkan tepat pada kasus. & Sebagian proses dijelaskan tepat, tetapi penerapan pada kasus belum lengkap. & Penjelasan proses eksekusi/penutupan keliru atau tidak relevan. \\ \\
Ketepatan pengelolaan komunikasi dan pemangku kepentingan & 5 & Strategi komunikasi dan pengelolaan pemangku kepentingan dijawab tepat sesuai kasus. & Jawaban cukup tepat, tetapi strategi belum lengkap. & Jawaban keliru atau tidak menjawab kasus. \\ \\
Ketepatan dokumentasi dan evaluasi proyek & 8 & Struktur dokumentasi, evaluasi, dan lessons learned proyek dijelaskan sistematis dan tepat. & Dokumentasi/evaluasi dijelaskan cukup tepat, tetapi belum sistematis. & Jawaban tidak menunjukkan kemampuan mendokumentasikan proyek. \\}

\begin{tblr}{width=\textwidth,colspec={|Q[l,m,3.2cm]|X[l,m]|},hlines,vlines,hspan=minimal,row{1-3}={bg=headgray},rowsep=1.5pt,colsep=3pt}
Jadwal Pelaksanaan: & Minggu 1--17 sesuai RPS. Setiap evaluasi memiliki RTM dan rubrik multi-kriteria. \\
Lain-Lain Yang Diperlukan: & LMS (lmsslc.polinema.ac.id), dokumen proyek kelompok, perangkat lunak sesuai mata kuliah, dan perangkat presentasi. \\
Pustaka: & Mengacu pustaka utama dan pendukung pada bagian RPS. \\
\end{tblr}
